\section*{Capítulo \ref{ch:b3:nervous-systems} --- Organização dos sistemas nervosos}

\begin{solution}{exo:b3:nervous-systems:1}
\omterm{def:b3:nervous-systems:cells}{Astrócitos}: tamponam potássio e glutamato, alimentam os neurônios,
induzem a \omterm{prop:b3:nervous-systems:barrier-budget}{barreira hematoencefálica}, regulam o fluxo sanguíneo local.
\omterm{def:b3:nervous-systems:cells}{Oligodendrócitos}: mielinizam axônios no sistema nervoso central. Células
de Schwann: mielinizam axônios periféricos (e guiam sua regeneração).
\omterm{def:b3:nervous-systems:cells}{Micróglia}: os \omterm{def:b3:innate-immunity:innate}{macrófagos} do encéfalo, que podam sinapses e removem
detritos. (As células ependimárias revestem os ventrículos e fazem o
\omterm{prop:b3:nervous-systems:barrier-budget}{líquido cefalorraquidiano}.)
\end{solution}

\begin{solution}{exo:b3:nervous-systems:2}
A coloração enegrecia por inteiro alguns poucos neurônios e deixava os
demais invisíveis, de modo que uma única célula podia ser vista
completa. Cajal concluiu que os neurônios são células separadas que se
contatam sem se fundir, com uma direção fixa de fluxo dos dendritos ao
axônio; Golgi concluiu que formavam uma rede contínua. Cajal estava
certo; o microscópio eletrônico mostrou a fenda sináptica em 1954.
\end{solution}

\begin{solution}{exo:b3:nervous-systems:3}
Uma batida estira o quadríceps; seus fusos musculares disparam; os
aferentes Ia entram na medula e fazem sinapse diretamente sobre os
neurônios motores do quadríceps, que disparam e contraem o músculo,
enquanto um \omterm{def:b3:nervous-systems:cells}{interneurônio} inibe os neurônios motores dos isquiotibiais.
Uma sinapse, axônios curtos a \qty{80}{m/s}: cerca de \qty{30}{ms}. Um
chute voluntário passa pelo encéfalo --- percepção, decisão,
planejamento motor, tratos descendentes, várias sinapses --- e leva
\qtyrange{200}{300}{ms}.
\end{solution}

\begin{solution}{exo:b3:nervous-systems:4}
Simpático (noradrenalina no alvo): coração mais rápido e mais forte,
pupila dilatada, motilidade e secreção intestinais reduzidas, vias
aéreas dilatadas. Parassimpático (acetilcolina): coração desacelerado,
pupila contraída, motilidade e secreção intestinais aumentadas, vias
aéreas contraídas.
\end{solution}

\begin{solution}{exo:b3:nervous-systems:5}
$\lambda = \sqrt{R_{m}d/4R_{i}} = \sqrt{10^{4}\times 10^{-4}/400}$ cm
$= \qty{0.05}{cm} = \qty{0.5}{mm}$. Depois de \qty{2}{mm}: $e^{-4} \approx
2\,\%$. Dobrar $\lambda$ exige quatro vezes o diâmetro,
\qty{4}{\micro m}.
\end{solution}

\begin{solution}{exo:b3:nervous-systems:6}
Sensorial: $12\times 6 = \qty{72}{m/s}$, $1.2/72 = \qty{17}{ms}$. Motor:
\qty{90}{m/s}, $1.1/90 = \qty{12}{ms}$. Duas sinapses, \qty{1}{ms}:
cerca de \qty{30}{ms}. Amielínicos: $1.1\sqrt{12} = \qty{3.8}{m/s}$ e
$1.1\sqrt{15} = \qty{4.3}{m/s}$: $315 + 258 + 1 \approx \qty{570}{ms}$
--- lento demais para salvar o pé.
\end{solution}

\begin{solution}{exo:b3:nervous-systems:7}
$2.4\times 10^{20}/8.6\times 10^{10} \approx 2.8\times 10^{9}$ ATP por
neurônio por segundo; a $2\times 10^{9}$ por disparo, cerca de
\qty{1.4}{Hz} em média. O código tem de ser esparso: poucos disparos,
com a informação carregada por quais neurônios disparam e quando, e não
por taxas altas em toda parte.
\end{solution}

\begin{solution}{exo:b3:nervous-systems:8}
Dois neurônios que se inibem mutuamente sem se cansar formam um
interruptor: o primeiro a disparar suprime o outro enquanto durar o
estímulo, e nunca há troca. Oscilar exige um processo lento que
enfraqueça o domínio do vencedor --- adaptação de seu disparo, depressão
de sua sinapse inibitória, ou um rebote do perdedor após a liberação ---
cuja constante de tempo fixa o período.
\end{solution}

\begin{solution}{exo:b3:nervous-systems:9}
A L-dopa é carregada através da barreira pelo transportador de
aminoácidos neutros grandes; a dopamina, polar e carregada, não é. Os
pacientes tomam L-dopa (com um inibidor da enzima periférica, para que
ela não seja convertida antes de atravessar), e a enzima do próprio
encéfalo faz dopamina a partir dela onde é necessária. Os \omterm{def:b3:adaptive-immunity:antibody}{anticorpos},
com \qty{150}{kDa}, não atravessam nem as junções oclusivas nem
transportador algum; a entrega exige um transportador de carona que
sequestre um receptor para transcitose, ou injeção no \omterm{prop:b3:nervous-systems:barrier-budget}{líquido
cefalorraquidiano}.
\end{solution}

\begin{solution}{exo:b3:nervous-systems:10}
A corrente de membrana por unidade de comprimento é a resistiva
$V/r_{m}$ mais a capacitiva $c_{m}\,\partial V/\partial t$; a
conservação de carga dá $\partial I/\partial x = -(V/r_{m} + c_{m}\,\partial V/\partial t)$ e, com $\partial V/\partial x = -r_{i}I$,
$\frac{1}{r_{i}}\,\partial^{2}V/
\partial x^{2} = V/r_{m} + c_{m}\,\partial V/\partial t$; multiplicando por $r_{m}$:
$\lambda^{2}\,\partial^{2}V/\partial x^{2} = V + \tau\,\partial
V/\partial t$. Com $\partial V/\partial t = 0$ retorna a equação do
teorema. $\lambda$ é um comprimento e $\tau$, um tempo, de modo que
$\lambda/\tau$ é uma velocidade --- a escala natural da rapidez com que
uma perturbação se espalha.
\end{solution}

\begin{solution}{exo:b3:nervous-systems:11}
$v \propto \sqrt{d}$: quatro vezes a velocidade exige dezesseis vezes o
diâmetro, \qty{8}{mm}. Área $\pi(4\,\text{mm})^{2} \approx \qty{50}{mm^2}$ contra $\pi\,(10\,\unit{\micro m})^{2} = \qty{314}{\micro m^2}$: uma razão de $1.6\times 10^{5}$. Um nervo de
fibras amielínicas rápidas é impossível --- uma única fibra dessas seria
tão grossa quanto o nervo ---, de modo que um animal com muitos canais
rápidos precisa de \omterm{def:b3:nervous-systems:cells}{mielina}, e é por isso que os vertebrados (e uns
poucos invertebrados, independentemente) a desenvolveram.
\end{solution}

\begin{solution}{exo:b3:nervous-systems:12}
$a = 0.4/30^{0.75} = 0.4/12.8 = 0.031$; para \qty{70000}{g}: $0.031\times
\num{70000}^{0.75} = 0.031\times 4300 \approx \qty{134}{g}$. Os
\qty{1350}{g} humanos são dez vezes a previsão. A massa encefálica segue
a massa corporal porque um corpo maior precisa de mais neurônios
sensoriais e motores, e porque os próprios neurônios crescem com o
tamanho do encéfalo; o que acompanha a cognição é o número de neurônios
do córtex, que depende de quão densamente eles são empacotados --- os
primatas empacotam mais por grama que os demais mamíferos, e os humanos
mais que todos.
\end{solution}

\begin{solution}{pb:b3:nervous-systems:1}
\textbf{1.} $\lambda = \sqrt{2\times 10^{4}\,d/400}$: para $d =
\qty{1e-4}{cm}$, $\sqrt{5\times 10^{-3}} = \qty{0.071}{cm} =
\qty{0.71}{mm}$; para \qty{4}{\micro m}, \qty{1.41}{mm}. $\tau = 2\times
10^{4}\times 10^{-6} = \qty{20}{ms}$.
\textbf{2.} $5\,e^{-0.3/0.71} = \qty{3.3}{mV}$; $5\,e^{-0.3/1.41} =
\qty{4.0}{mV}$.
\textbf{3.} As entradas distais chegam atenuadas, de modo que o lugar
onde a sinapse cai fixa seu peso; um potencial decai em cerca de $\tau$,
de modo que duas entradas se somam apenas se caírem dentro de uns
\qty{20}{ms} uma da outra --- o neurônio é um detector de coincidências
com uma janela de \qty{20}{ms}.
\textbf{4.} Mielínico, \qty{60}{m/s}; amielínico, $1.1\sqrt{10} =
\qty{3.5}{m/s}$; para chegar a \qty{60}{m/s} sem \omterm{def:b3:nervous-systems:cells}{mielina}, $d = (60/1.1)^{2}
\approx \qty{3000}{\micro m}$: três milímetros.
\textbf{5.} $\pi(5)^{2} = \qty{79}{\micro m^2}$ contra $\pi(1500)^{2} =
7\times 10^{6}\,\unit{\micro m^2}$: cerca de \num{90000} axônios
mielínicos no espaço de um só.
\textbf{6.} \qty{2}{cm} a \qty{60}{m/s} são \qty{0.3}{ms}; a
\qty{3.5}{m/s}, \qty{5.7}{ms}: uns \qty{5}{ms} de atraso extra. Pior, a
membrana desnuda tem capacitância grande e poucos canais, de modo que a
corrente vinda do último nodo intacto pode não levá-la ao limiar:
bloqueio de condução.
\textbf{7.} Sensorial, $1.2/60 = \qty{20}{ms}$; duas sinapses,
\qty{1}{ms}; motor, $1.1/90 = \qty{12}{ms}$: cerca de \qty{33}{ms} (mais
um milissegundo na junção neuromuscular).
\textbf{8.} Até o encéfalo: $0.6/60 = \qty{10}{ms}$ mais duas sinapses,
\qty{11}{ms} depois de o sinal entrar na medula aos \qty{20}{ms}: o
sinal chega ao córtex mais ou menos quando o pé se move; a experiência
consciente da dor exige mais algumas centenas de milissegundos de
processamento, de modo que o pé é retirado antes de a dor ser sentida.
\textbf{9.} $1.1\sqrt{1} = \qty{1.1}{m/s}$: $1.8/1.1 \approx
\qty{1.6}{s}$. A primeira dor, aguda, chega por fibras mielínicas finas
em dezenas de milissegundos; a segunda dor, surda e em queimação, pelas
fibras C amielínicas um segundo ou mais depois.
\textbf{10.} Sensorial, $3.2/60 = \qty{53}{ms}$; motor, $3.1/90 =
\qty{34}{ms}$; sinapses, \qty{1}{ms}: cerca de \qty{90}{ms}. Animais
compridos têm reflexos lentos; eles compensam com axônios mais grossos e
delegando mais aos circuitos locais.
\textbf{11.} $0.8/80 = \qty{10}{ms}$ em cada sentido, \qty{20}{ms};
sinapse, \qty{0.5}{ms}; junção, \qty{1}{ms}; força, \qty{5}{ms}:
\qty{26.5}{ms}, com a própria transdução do fuso e a condução dentro do
músculo completando o resto dos \qty{30}{ms}.
\textbf{12.} Os axônios amielínicos conduzem a um ou dois metros por
segundo, de modo que alças que levam \qty{30}{ms} num adulto levam
centenas, com temporização ruim; a mielinização eleva a velocidade de
dez a cinquenta vezes e torna a temporização precisa, que é o que andar,
agarrar e falar exigem.
\textbf{13.} $20/50\,000 = 4\times 10^{-4}$ mol/s $= 2.4\times 10^{20}$
ATP por segundo; $2.8\times 10^{9}$ por neurônio por segundo.
\textbf{14.} $4\times 10^{8} + 7000\times 2\times 10^{4} = 4\times 10^{8}
+ 1.4\times 10^{8} = 5.4\times 10^{8}$ ATP.
\textbf{15.} Tudo em disparos: $2.8\times 10^{9}/5.4\times 10^{8} \approx
\qty{5}{Hz}$; metade: cerca de \qty{2.6}{Hz}.
\textbf{16.} A \qty{1}{Hz}, $5.4\times 10^{8}/2.8\times 10^{9} \approx
\qty{20}{\%}$ do orçamento: o código é esparso --- a maioria dos
neurônios em silêncio a maior parte do tempo, com a informação em quais
poucos disparam e quando.
\textbf{17.} $8.6\times 10^{10}\times 50\times 5.4\times 10^{8} = 2.3
\times 10^{21}$ ATP/s $= 3.9\times 10^{-3}$ mol/s $\approx \qty{190}{W}$
--- o dobro de todo o corpo em repouso.
\textbf{18.} $5\times 16 = \qty{80}{kJ/h} = \qty{22}{W}$: exatamente o
suficiente, sem reserva. Quando o abastecimento para, as bombas param em
segundos, os gradientes se dissipam e a consciência se perde em cerca de
dez segundos.
\textbf{19.} $a = 0.4/12.8 = 0.031$; \qty{70}{kg}: $0.031\times 4300 =
\qty{134}{g}$; \qty{5000}{kg}: $0.031\times 1.06\times 10^{5} \approx
\qty{3300}{g}$.
\textbf{20.} Humano, $1350/134 \approx 10$; elefante, $4800/3300 \approx
1.5$.
\textbf{21.} No \omterm{def:b3:nervous-systems:plan}{cerebelo} --- \qty{98}{\%} deles ---, coordenando os
quarenta mil músculos da tromba. O que acompanha a cognição é o número
de neurônios corticais, e não o total; o \omterm{def:b3:nervous-systems:plan}{cerebelo} acompanha a tarefa
motora de comandar um corpo enorme.
\textbf{22.} Os axônios e dendritos se alongam com o encéfalo, de modo
que cada neurônio ocupa mais volume e o número de neurônios cresce mais
devagar que a massa. Nos primatas o tamanho dos neurônios permanece
quase constante enquanto o encéfalo cresce, de modo que o número de
neurônios cresce quase em proporção à massa, e um encéfalo de primata de
dado peso abriga várias vezes os neurônios de um encéfalo de roedor
desse peso.
\textbf{23.} Um braço consegue coordenar suas próprias ventosas e
articulações, alcançar, agarrar, passar comida ao longo de si mesmo e
evitar agarrar a própria pele, por circuitos locais; não consegue ver,
escolher um alvo nem aprender uma discriminação. Teste: separe o cordão
nervoso do braço do encéfalo e toque o braço com comida --- ele agarra e
a passa em direção a onde estaria a boca; um braço amputado faz o mesmo
por uma hora.
\textbf{24.} Cerca de $23$ sinapses por neurônio contra $7000$:
trezentas vezes menos. O mapa completo diz qual neurônio pode influenciar
qual, mas não a força nem o sinal de cada conexão, nem que
neuromoduladores as alteram, nem a dinâmica --- de modo que, mesmo para o
verme, o diagrama da fiação prevê o comportamento apenas em parte.
\textbf{25.} $\lambda = \qty{0.71}{mm}$ para o dendrito de
\qty{1}{\micro m}; latência de retirada de cerca de \qty{33}{ms};
frequência média de disparo de cerca de \qty{2.6}{Hz} com metade do
orçamento em disparos.
\end{solution}
